\paragraph{Distinguish collecting a claim from assessing whether it is overdue.}
The same receivable balance can imply different collection problems at different
dates. Consider two explicitly assumed claims: USD~60 due on 31~January 2027
and USD~40 due on 28~February. A receipt of USD~20 against the first claim on
15~January leaves USD~40 against that claim. The selected synthetic policy treats
a remaining amount as overdue only after its due date. Thus neither claim is
overdue at the end of January. A further receipt of USD~10 against the second
claim in February leaves aggregate receivables of USD~70. At 28~February, USD~40
is overdue under this policy; the other USD~30 becomes overdue on the next day.
These classifications describe the assumed payment terms. They do not measure
an allowance or establish a right to write off either claim.

\texttt{DeclaredCollectionClock/1} keeps a fixed declared population of claims,
with separate memory for each due date, remaining amount and last assessed
amount. The opening claims must sum exactly to the selected receivable account.
A collection identifies one claim, cannot exceed its recorded remainder, and
posts debit cash and credit receivables. It changes that claim's remainder while
preserving the previous assessment. An assessment instead reads the remainders
at its declared date, stores the resulting overdue amounts and makes no journal
entry or cash movement. Unknown required event inputs retain their reasons and
stop that command; an explicitly false applicability finding refuses the method.

Assessment windows have an inclusive start and an exclusive end. The assessment
occurs on the final included date, and the saved cursor must equal that window's
start before it can advance to the end. Windows must be positive and contiguous.
This makes a repeated, overlapping or skipped assessment detectable. Within one
date, supplied event order is retained, but an assessment seals that date against
a later collection: a payment intended to affect that observation must precede
it. The journal engine also rejects events earlier than the current book date
or outside the current fiscal period.

A report can advance its observation date without supplying the missing
assessment. It then identifies the last assessment as historical, or records
that none exists. Provisional period carry preserves unpaid claims, due dates
and assessment memory; it cannot reset the cursor or turn year-end into a new
assessment. In the example, a December assessment records overdue amounts of
USD~40 and USD~30. After carrying those balances, collecting the first USD~40
in January 2028 leaves USD~30 outstanding. A declared January assessment then
records overdue amounts of zero and USD~30. Cash has risen from USD~100 to
USD~170, receivables have fallen from USD~100 to USD~30, and total assets and
equity remain USD~200 throughout.

The existing command language can author these events, reports and carries
through its hashed definitions. The guarded execution compares exact method
definitions, so renaming a registry entry does not change its meaning. A foreign
method that writes claim memory or posts to the receivable account refuses,
including offsetting postings with no net balance change. Only the successful
prefix remains available after refusal. The checked fragment supports eight
declared opening claims and sixteen assessment windows within the existing work
limits. New claim insertion, revisions to payment terms, complete invoice
coverage, impairment, tax and source admission remain distinct requirements.

